\section{三位皇帝与三种立国条件}

公元二二〇年正月，曹操去世。曹丕接过魏王位，也接过已经由曹氏掌握的朝廷军政机构。十月，献帝禅位，曹丕即皇帝位，改国号为魏。受禅的诏册把这一变化写成可以由礼仪完成的王朝交接；但礼仪之所以能够举行，是因为魏王国先已掌握北方的军队、官署和发出诏令的中枢。汉朝皇统至此不再是实际的统治中心。要紧的不是把献帝的意愿简化为一个可以揣测的内心问题，而是看见一套仍用汉廷名义运转的机构，如何已经为曹氏的继承服务。

魏随即把都城置于洛阳。许县更接近汉末的战时行营，洛阳则把尚书机构、宫廷警卫、北方仓储与中原道路重新接在一起。这个选择说明魏的资源基础在哪里：它拥有华北较广阔的征发区域、既有官僚与军事网络，也能够把旧都的政治象征转化为新朝的行政枢纽。不过，都城并不能自行解决边界问题。魏控制的北方财富须经由淮汉、汉水上游和山地关隘才能成为对吴、蜀的压力；越过这些地带时，舟师、季节、道路和地方守军都会把纸面上的优势拆成具体而有限的行动。

\subsection{汉统的名义与益州的边界}

曹丕受禅后，刘备没有承认魏的皇帝地位。次年四月，他在成都即皇帝位，建元章武，以汉室继承者自居。成都朝廷的称帝，并不只是同魏争一个称号：它把原来的汉中王政权改置为可任命官员、分配爵位、征发兵役的皇帝体制，也向益州官民说明此处并非接受魏命的地方军府。魏以禅代取得汉的名义，蜀汉则把汉室延续作为自己的合法性语言；两种语言都指向汉朝留下的制度和记忆，却不能使任何一方立刻取得天下的服从。

蜀汉的可用根基比其名义更有边界。益州的山川、物产和成都的官府给它以一处稳固腹地，汉中则是面对魏的山口屏障；但它的主要资源仍在益州和汉中，已不再握有荆州。关羽败亡以后，荆州的失去既截断东进通道，也使成都朝廷面对领土与集团内部的压力。刘备在二二一年集结益州军东下，正以夺回荆州和为关羽复仇为理由。这里不宜把一场战争归结为一种单纯动机：史料能够确知的是，失地、盟友关系的破裂与集团的诉求同时构成了出兵的条件。

孙权一面要抵御蜀军，一面要应对北方。二二二年，曹丕册封孙权为吴王，并以荆州牧名义和交出质子的要求，试图把长江下游政权放进魏的册封秩序。孙权暂时接受封号，可以牵制正向东而来的蜀军；这并不表示江东的军政资源已经交给洛阳。册封是外交中的名义安排，不是对长江水面、州郡吏员和地方兵力的自动接收。魏想利用孙刘相争扩大影响，孙权则利用魏的名号延缓眼前的两线压力，三方的合法性语言因而始终伴随着相互防备。

\subsection{长江没有替任何一方决定胜负}

同年，陆逊统率吴军在夷陵、猇亭一线击破刘备，蜀军沿江退向白帝。战役的胜负与大致地点可以相互印证；火攻如何展开、营垒究竟延展多远、伤亡有多少，以及陆逊如何取得全军信任，传世叙事并不提供同样牢靠的答案。较稳的解释不必借助这些戏剧化细节：蜀军深入荆州，补给线拉长，连营难以灵活转移；吴军则依托长江中游的水陆节点，等待对手疲敝。夷陵使蜀汉失去立即夺回荆州的可能，也让孙吴守住了上游屏障。

公元二二三年四月，刘备在永安去世，刘禅继位，诸葛亮受遗诏辅政。皇帝更替发生在夷陵败退之后，蜀汉既要防备孙吴，又要防止益州内部因继承而离散。辅政体制使军政命令能够延续，却不能抹去小国资源有限的事实。邓芝同年出使孙吴，双方恢复通好，孙刘再度形成共同抗魏的外交框架。它稳定了长江上游的战线，使蜀汉可以转向南中与汉中；但荆州归属造成的裂痕并未因此消失。盟约的具体条款和使节言辞难以完全复原，恢复通好这一结果则比那些言辞更能说明两国的选择：各自都不能承受让魏独占北方资源后的两面压迫。

魏的限制也在江面上显露出来。二二四年，曹丕到广陵观察南征，渡江条件终因冰冻、舟师和前进基地不足而不具备。史书并不能据此给出可靠的兵员总数，曹丕意在试探、决战还是兼有二者也不宜断言；撤军本身已经说明，控制北方仓储和陆路并不等于拥有可跨越长江的水军能力。长江不是天然替孙吴保证安全的墙，却迫使进攻者把季节、水文、船只与补给纳入同一项计算。

\subsection{把腹地变成国家能力}

夷陵之后，三国都在把已有的资源重新编入可以持续动员的秩序。二二五年，诸葛亮南征，平定益州郡、永昌等地的叛乱，重新调整蜀汉与南中地方集团的关系。成都需要南中的铜铁、赋役、人力和通往西南的道路；这场行动因此既是军事平乱，也关系到小国能否负担边防。所谓孟获及“七擒”的故事，主要依赖较晚、层次复杂的材料，孟获的身份、次数、行军路线和具体处置都不能按通行叙事照录。南征的发生及其大致结果较为可靠，却不等于地方社会从此被完全整合。地方首领、夷帅、郡县官和成都朝廷之间的关系，仍只能从有限的叙事中窥见。

二二六年曹丕去世，曹叡继位。魏的继承能够迅速完成，显示洛阳中枢和北方国家机器已有相当的连续性；与此同时，它也必须同时面对汉中、淮南和江淮方向的战线。蜀汉在南中大体安定、又与孙吴复盟后，于二二七年把军政准备移至汉中。诸葛亮驻守汉中并上《出师表》，把北伐说成维系先帝遗志与国家存续的行动。《出师表》的传本虽为后人理解这一选择留下重要文字，但其成文语境和是否完整保留了当时的决策讨论，仍须保留疑问。可以确定的是，北伐必须从汉中越过山地，补给距离和关隘会持续限定蜀军的规模与停留时间。

二二八年，蜀军出祁山，南安、天水、安定一度响应；魏调张郃等西援，随后马谡失守街亭，诸葛亮撤回汉中。数郡为何响应、倒向的范围有多大，不能从“震动”一词推出全面归附；这一插曲至少表明，边郡并非始终以同一种方式看待魏蜀。街亭的失守则使通向陇右的要道无法稳固，蜀军不能维持前出部署。后世常把败因凝为一名将领的过错，但高地、水源、道路、补给和指挥的相互牵制，比一个道德化结论更接近战争的约束。

同在二二八年，孟达在新城的叛乱被司马懿迅速平定，魏保住了汉水上游的节点；陆逊又在石亭击败曹休，魏在江淮方向的攻势受挫。前一事件所涉孟达与蜀汉的书信，多见于裴松之注所引材料，真伪和行军日程不能不加辨析；后一战的诈降情节亦主要见于传记性铺陈。把这些细节暂置一旁，魏在西部的快速机动与孙吴在江淮水陆节点上的应对，仍显示三国不是静止的三块疆域，而是数条补给线相互试探的竞争。

二二九年，诸葛亮取得武都、阴平二郡，改善汉中西北的屏障，却未能由此穿透魏的渭水防线。郡界、守令转换及实际控制到达何处，须逐县核验，不能将两郡名义写成整个陇右的转向。同年四月，孙权在武昌即皇帝位，建元黄龙；此后又把都城迁回建业。魏的册封和质子要求已不足以容纳江东的自主统治，蜀汉也早已称帝，孙吴遂以自己的年号、官爵和都城提出独立的皇帝中心。建业位于长江下游的交通与资源腹地，武昌仍有面向上游的军事作用；迁都不是把一个城市简单替换为另一个城市，而是调整江东官僚、航运和边防的重心。

\subsection{史书的三国与可见的三国}

这一段年序的主干，首先依靠《三国志》：其中魏的纪年居于全书的框架，蜀汉和孙吴的材料分别进入各自的记述，因此同一事件在篇幅、称谓和视角上并不对称。裴松之的注保存了《魏略》等已佚材料和许多异说，尤其有助于看见孟达书信、马谡处置等叙事的来源层次；注引并不因其丰富就与陈寿正文具有同等证明力。《后汉书》〈献帝纪〉可用来校看禅代一端，《华阳国志》补出巴蜀与南中视角，《资治通鉴》则把三方事件编为连续年序。它们都是必要的交叉材料，却各有后出编纂、区域视角或正统排序的取舍，不能把相合的叙事误当作互不相关的现场见证。

考古与文书让这幅图景多出史书之外的尺度。汉魏洛阳城遗址可以检验旧都作为政治空间的延续，南京地区的城址材料可帮助追问建业长期演变；两者都不能直接复原二二〇年或二二九年的每一座宫城、人口数和迁徙路线。走马楼吴简保留了长沙郡临湘县的户籍、田租和吏役痕迹，说明孙吴的统治必须落到簿籍和基层行政上；它毕竟是特定地点与机构的文书，不能替孙吴全境发言。正因如此，魏、蜀汉、孙吴的建立不应被写成三次仪式性的登场：受禅、称帝与迁都确立了三个皇帝中心，而它们能否延续，仍取决于各自能征发什么、能沿何种道路运送、又在哪些河流与山口前停住。
